%!TEX root = ../template.tex
\typeout{NT FILE abstract-pt.tex}%


O Grande Panorama de Lisboa, painel de azulejos de 23 metros atribuído a
Gabriel del Barco (c.1700), é o registo pictórico mais completo de Lisboa
anterior ao terramoto de 1755, catalogado no encerrado Museu Nacional do
Azulejo (MNAz). Apesar da escala, uma legenda estática
torna os seus 150+ edifícios ininterpretáveis sem conhecimento
especializado. Este relatório propõe a TileStories, framework de
realidade aumentada orientada por configuração, que sobrepõe conteúdo
adaptado ao visitante sobre uma parede patrimonial em torno de um eixo
de estado generalizado --- no Panorama, uma linha temporal de quatro
épocas (antes de 1755, o terramoto, a reconstrução pombalina e a
atualidade) --- com personalização inicial que adapta conteúdo e
percursos a perfis de visitante. O contributo técnico central é
configurar e avaliar o Immersal, serviço de posicionamento visual, para
paredes de azulejo extensas e contínuas com múltiplos
pontos de interesse, numa abordagem iterativa de Design Science
Research. O principal risco técnico já está demonstrado: uma integração
funcional do Immersal localiza-se numa parede real em dois a três
segundos. A generalização é testada em três casos
distintos: o Panorama; o Chafariz Velho (Paço de Arcos),
mais exigente e superfície principal de desenvolvimento durante o
encerramento do MNAz; e um mural contemporâneo de fauna selvagem sem
estatuto equivalente. Espera-se que a tese contribua com uma
configuração de rastreamento multi-zona validada, um mecanismo de
personalização do visitante, um mecanismo de visualização de estado
generalizado, e uma avaliação in situ de envolvimento, usabilidade e
aprendizagem com visitantes reais.
 
\keywords{Realidade Aumentada, Património Cultural, Personalização do
Visitante, Aplicações Móveis, Visualização Histórica, Design Science
Research, Grande Panorama de Lisboa}

